
El modelo DEB formulado en el \cref{chap:cap3} contiene parámetros de dos naturalezas: (i) \emph{semillas}, derivados directamente de la literatura bioenergética para \textit{Hermetia illucens}, que se mantienen fijos por su fundamento experimental; y (ii) \emph{libres}, que dependen de condiciones locales del reactor (geometría, sustrato, densidad de siembra) y deben ajustarse mediante optimización contra las series de sensores. Esta sección detalla la estrategia de calibración en dos fases secuenciales, culminando en la estimación restringida de los parámetros libres.

\subsection{Fase A: semillas de parámetros a partir de Eriksen (2022)}
\label{subsec:cap4_fase_a}

Los parámetros semilla se extraen del modelo DEB individual de \textcite{eriksenDynamicModellingFeed2022}, escalados a nivel poblacional mediante la densidad larval inicial $B_0$. Estos valores se consideran \emph{a priori} biofísicamente plausibles y no se alteran durante la calibración local:

\begin{table}[htbp]
    \centering
    \caption{Parámetros semilla del modelo DEB, fijos durante la calibración.}
    \label{tab:cap4_parametros_semilla}
    \begin{tabular}{@{}llll@{}}
        \toprule
        \textbf{Símbolo} & \textbf{Valor} & \textbf{Unidad} & \textbf{Significado biológico} \\
        \midrule
        $\kappa$ & \num{0.60} & adimensional & Fracción de energía asimilada al soma \\
        $m$ & \num{2.1e-4} & min$^{-1}$ & Tasa específica de mantenimiento \\
        $a_{\max}$ & \num{1.25e-4} & g C$^{1/3}$/min & Coeficiente de asimilación máxima \\
        $Y_B$ & \num{0.80} & adimensional & Eficiencia de conversión a biomasa \\
        $Y_L$ & \num{0.75} & adimensional & Eficiencia de conversión a lípidos \\
        \bottomrule
    \end{tabular}
\end{table}

La fracción $\kappa = \num{0.60}$ implica que el \num{60}\% de la energía asimilada se destina a mantenimiento y crecimiento somático, mientras que el \num{40}\% restante se reserva para reproducción (en larvas, acumulación de lípidos de reserva). La tasa de mantenimiento $m$ refleja el costo metabólico basal por unidad de biomasa estructural, y $a_{\max}$ codifica la capacidad máxima de ingestión por unidad de superficie corporal, asumiendo una isometría volumen-superficie típica de insectos \parencite{eriksenDynamicModellingFeed2022}.

\subsection{Fase B: ajuste local con series de sensores NDIR}
\label{subsec:cap4_fase_b}

Los parámetros libres capturan la variabilidad operativa entre réplicas y se ajustan minimizando la función de pérdida \eqref{eq:cap4_funcion_perdida} sobre las ventanas de confrontación modelo-sensor. El vector de parámetros libres es:
\begin{equation}
    \label{eq:cap4_parametros_libres}
    \boldsymbol{\theta}_{\mathrm{libre}} =
    \bigl[\,\tau_h,\; \beta_0,\; \rho_{\mathrm{conv}},\; T_A,\; T_{\mathrm{opt}},\; T_H\,\bigr]^\intercal
    \in \mathbb{R}^{6}.
\end{equation}

Cada componente tiene una interpretación física directa:
\begin{itemize}
    \item $\tau_h$: tiempo de manipulación del sustrato (min), que retrasa la disponibilidad de alimento asimilable.
    \item $\beta_0$: coeficiente de saturación del kernel de Holling II (g C), inversamente relacionado con la afinidad larval por el sustrato.
    \item $\rho_{\mathrm{conv}}$: factor de conversión de biomasa húmeda a carbono estructural (g C/g biomasa húmeda), que escala el pseudo-observable $\hat{B}(t)$ del GP al dominio del DEB.
    \item $T_A$: energía de activación de Arrhenius (K), que modula la tasa metabólica a bajas temperaturas.
    \item $T_{\mathrm{opt}}$: temperatura óptima del kernel de Sharpe-Schoolfield (°C).
    \item $T_H$: temperatura de inactivación de alta temperatura (°C).
\end{itemize}

La calibración se ejecuta sobre las $K$ ventanas de \num{30} min del ciclo experimental, excluyendo aquellas marcadas con $\mathcal{A}(t_k) = 1$ (alerta de anaerobiosis, véase \cref{subsec:cap4_tratamiento_ch4}). El conjunto de entrenamiento para la optimización contiene típicamente $K_{\mathrm{eff}} \approx \num{40}$ ventanas por ciclo de \num{14} días.

\subsection{Restricciones de caja y método de optimización}
\label{subsec:cap4_restricciones}

Para preservar la interpretabilidad biológica y evitar mínimos locales no físicos, cada parámetro libre se acota mediante restricciones de caja inferidas del conocimiento zootécnico de \textit{H. illucens}:

\begin{table}[htbp]
    \centering
    \caption{Restricciones de caja para los parámetros libres del modelo DEB.}
    \label{tab:cap4_restricciones_caja}
    \begin{tabular}{@{}lccc@{}}
        \toprule
        \textbf{Parámetro} & \textbf{Límite inferior} & \textbf{Límite superior} & \textbf{Justificación} \\
        \midrule
        $\tau_h$ & \num{10} min & \num{120} min & Tiempo de fragmentación mecánica \\
        $\beta_0$ & \num{0.1} g C & \num{5.0} g C & Rango de densidad larval típica \\
        $\rho_{\mathrm{conv}}$ & \num{0.05} & \num{0.15} & Fracción de carbono en biomasa húmeda \\
        $T_A$ & \num{2000} K & \num{8000} K & Energía de activación de enzimas digestivas \\
        $T_{\mathrm{opt}}$ & \num{25} °C & \num{35} °C & Rango de confort térmico larval \\
        $T_H$ & \num{38} °C & \num{45} °C & Temperatura letal de inactivación \\
        \bottomrule
    \end{tabular}
\end{table}

El problema de optimización se plantea como:
\begin{equation}
    \label{eq:cap4_optimizacion}
    \hat{\boldsymbol{\theta}}_{\mathrm{libre}} =
    \arg\min_{\boldsymbol{\theta} \in \Theta}
    \mathcal{L}(\boldsymbol{\theta}),
    \qquad
    \Theta = \prod_{d=1}^{6}\,[\theta_d^{\min},\, \theta_d^{\max}],
\end{equation}
donde $\mathcal{L}(\boldsymbol{\theta})$ es la función de pérdida definida en la \cref{eq:cap4_funcion_perdida}.

Dado que el gradiente analítico de $\mathcal{L}$ respecto a $\boldsymbol{\theta}$ es impracticable —la función encadena integración numérica de EDO, interpolación del GP y derivadas de series temporales— se recurre a un optimizador de gradiente aproximado. Se emplea el algoritmo L-BFGS-B (Limited-memory Broyden-Fletcher-Goldfarb-Shanno con restricciones de caja), disponible en la biblioteca SciPy de Python. El gradiente se aproxima mediante diferencias finitas centradas de paso adaptativo:
\begin{equation}
    \label{eq:cap4_gradiente_df}
    \frac{\partial \mathcal{L}}{\partial \theta_d} \approx
    \frac{\mathcal{L}(\theta_d + h_d) - \mathcal{L}(\theta_d - h_d)}{2h_d},
    \qquad
    h_d = \num{1e-4}\,|\theta_d| + \num{1e-6},
\end{equation}
evitando la dependencia de frameworks de diferenciación automática (autodiff) que complicarían el despliegue en el microcontrolador ESP32 del nodo edge \parencite{weichertReviewMachineLearning2019}.

La optimización se inicializa en el centroide del hiperrectángulo $\Theta$ y se ejecuta hasta que la norma del gradiente proyectado cae por debajo de $\num{1e-4}$ o el cambio relativo en $\mathcal{L}$ es inferior a $\num{1e-6}$ durante \num{5} iteraciones consecutivas. En la práctica, la convergencia requiere entre \num{80} y \num{150} evaluaciones de la función de pérdida, lo que se completa en menos de \num{10} min en una estación de trabajo estándar, viabilizando su ejecución offline entre ciclos experimentales.

El resultado de la calibración es un vector $\hat{\boldsymbol{\theta}}_{\mathrm{libre}}$ específico por réplica, que captura la huella operativa del reactor bajo condiciones locales. La variabilidad entre réplicas de $\hat{\boldsymbol{\theta}}_{\mathrm{libre}}$ informa sobre la reproducibilidad del proceso y la robustez del sistema ante la heterogeneidad del sustrato.
